\documentclass[xcolor=dvipsnames]{beamer}
%\usepackage[utf8]{inputenc}
%\usepackage{xcolor}
\usepackage{graphicx}
\usepackage{MnSymbol}
\usepackage{tikz}
\usepackage{media9}
\usetikzlibrary{arrows,shapes}
\usepackage{stmaryrd}
\usepackage{colortbl}
\usepackage{caption}
\usepackage{comment}
%\usepackage[utf8]{inputenc}
\usepackage[british,UKenglish,USenglish,english,american]{babel}
%\usepackage[utf8]{vietnam}
\usepackage{pdfpages}
\usepackage{listings}
\usepackage{color}
\usepackage{booktabs}
\usepackage{soul}
\usepackage[normalem]{ulem}

%%%%%%%%%%%%%%%%%%%%%%%%%%%%%%%%%%%%%%%%%%%%%%%%%%%%%

%%%%%%%%%%%%%%%%%%%%%%%%%%%%%%%%%%%%%%%%%%%%%%%%%%%%%
\usepackage{tcolorbox}
\usepackage{lipsum}


%%%%%%%%%%%%%%%%%%%%%%%%%%%%%%%%%%%%%%%%%%%%%%%%%%%%%

\usepackage{pgf}
\usepackage{etex}
\usepackage{tikz,pgfplots}

\tikzstyle{every picture}+=[remember picture]
% By default all math in TikZ nodes are set in inline mode. Change this to
% displaystyle so that we don't get small fractions.
\everymath{\displaystyle}


\usetheme{Antibes}
%\usetheme{Madrid}
%\usecolortheme[named=Maroon]{structure}
\usecolortheme{dolphin}
\usefonttheme{professionalfonts}
\useoutertheme{infolines}
\useinnertheme{circles}

\newtheorem*{bem}{Bemerkung}

\usepackage{tikz}


%%%%%%%%%%%%%%%%%%%%%%%%%%%%%%%%%%%%%%%%%%%%%%%%%



%%%%%%%%%%%%%%%%%%%%%%%%%%%%%%%%%%%%%%%%%%%%%%%%%
\usepackage{listings}
\usepackage{color}

\definecolor{dkgreen}{rgb}{0,0.6,0}
\definecolor{gray}{rgb}{0.5,0.5,0.5}
\definecolor{mauve}{rgb}{0.58,0,0.82}

\lstset{frame=tb,
  language=Java,
  aboveskip=2mm,
  belowskip=12mm,
  showstringspaces=false,
  columns=flexible,
  basicstyle={\small\ttfamily},
  numbers=none,
  numberstyle=\tiny\color{gray},
  keywordstyle=\color{blue},
  commentstyle=\color{dkgreen},
  stringstyle=\color{mauve},
  breaklines=true,
  breakatwhitespace=true,
  tabsize=2
}
%%%%%%%%%%%%%%%%%%%%%%%%%%%%%%%%%%%%%%%%%%%%%%%%%

\title[International University - HCMC]{\includegraphics[height=1.8cm]{images/logo-iu.pdf}\\INTERNATIONAL UNIVERSITY - HCMC \\}
\author[DOAN Mau-Hien]{ \textbf{Title of Presentation}}
\institute{}
\logo{\includegraphics[height=1.5cm]{images/vnulogo.pdf}}

\logo{\pgfimage[height=0.8cm]{images/footweb.pdf}}


\begin{document}

\begin{frame}
  \titlepage
\end{frame}

\begin{frame}
\frametitle{\textbf{Outline:}}
\tableofcontents
\end{frame}

		
  
    \section{Section 1}  
     \begin{frame}{Title frame 1}
            \begin{block}{Block blue}
    		\begin{itemize}
    		\item Item 1
    		\item Item 2
    		\item Item 3
  			\end{itemize}
        \end{block}  
     \end{frame}
          
   \section{Section 2}  
     \begin{frame}{Title Frame 2}

\[\ support(X \to Y) = p(X \cup Y) = \frac{{n(X \cup Y)}}{N}\]
     \end{frame}
     
            
  
   
%*********************************************************************************************
\begin{frame}\frametitle{}
  \centering
    \includegraphics[scale=0.3]{images/thankyou}
    
\end{frame}

%**********************************************************************************************
%\begin{frame}\frametitle{}
%  \centering
%    \includemedia[
%  width=0.6\linewidth,height=0.45\linewidth,
%  activate=pageopen,
%  flashvars={
%    modestbranding=1 % no YT logo in control bar
%   &autohide=1       % controlbar autohide
%   &showinfo=0       % no title and other info before start
%  }
%]{}{http://www.youtube.com/v/dISaXUlilkU?rel=0}   % Flash file 
      
%\end{frame}
%%%%%%%%%%%%%%%%%%%%%%%%%%%%%%%%%%%%%%%%%%%%%%%%%%%
\end{document}
